% related_work.tex — Related Work (clustered by theme)
\section{Related Work}
\label{sec:related}

Rather than enumerate individual papers, we organize prior work into thematic clusters and highlight the evolution of ideas, recurring gaps, and occasional contradictions in the literature. Table~\ref{tab:literature} (see Section~\ref{sec:results}) provides a per-paper comparison.

\subsection{Object Detection for Accessibility}

The shift from hand-crafted features (HOG-SVM cascades, Haar classifiers) to deep convolutional detectors in the mid-2010s transformed what was possible for assistive vision. Redmon et al.'s YOLO family~\cite{redmon2016yolo} demonstrated that real-time, single-pass detection was achievable on commodity GPUs, and subsequent iterations~\cite{bochkovskiy2020yolov4, jocher2023yolov8} pushed accuracy while shrinking model size. At the same time, Liu et al.~\cite{liu2016ssd} showed that multi-scale detection via SSD could handle the wide size range of indoor obstacles, while Carion et al.~\cite{carion2020detr} replaced hand-designed anchors and NMS with a transformer-based end-to-end pipeline.

Several groups have applied these detectors specifically to blind-user scenarios. The DETR-based approach of Tran et al.~\cite{enhancing_detr_assistive2024} achieved strong mAP on curated indoor datasets but reported inference times of roughly 120\,ms on a desktop GPU---acceptable for research but tight for a real-time spoken interface that must also run depth, OCR, and TTS within the same 500\,ms budget. Zhang et al.~\cite{blind_aided_cascading2024} proposed a cascading feature pyramid with lightweight dual-path downsampling that reduced FLOPs by 38\%, though their evaluation was limited to synthetic corridor environments. More recently, transformer-based multimodal detection~\cite{transformer_multimodal_detection2024} has shown promise for fusing visual and linguistic cues, and assistive object recognition approaches~\cite{assistive_object_recognition2020} have explored personalized models.

A recurring gap across these studies is the \emph{absence of latency-aware evaluation}: most report mAP or F1 but not end-to-end time from frame capture to spoken alert. Our VVA system addresses this by enforcing per-stage timeouts (100\,ms detection, 300\,ms full pipeline) and providing graceful fallbacks when the primary detector is overloaded.

\subsection{Monocular Depth Estimation}

Ranftl et al.~\cite{ranftl2020midas} introduced MiDaS, a family of monocular depth networks trained on mixed datasets for zero-shot cross-domain transfer. The ``small'' variant (MiDaS v2.1, 256${\times}$256 input, ONNX export) has become a de facto baseline for assistive depth because it runs in under 30\,ms on a mid-range GPU. Kim et al.~\cite{obstacle_avoidance_uav_depth2022} applied fast monocular depth on UAVs and confirmed that metric accuracy degrades in textureless regions---a finding we reproduce in VVA's indoor tests. We mitigate this by pairing depth with edge-aware segmentation: when depth confidence is low, the segmenter's gradient-based boundary score provides an independent hazard signal.

\subsection{OCR and Document Accessibility}

Optical character recognition for assistive reading has advanced on two fronts: engine quality and preprocessing robustness. Smith's Tesseract~\cite{smith2007tesseract} remains the open-source stalwart, but EasyOCR~\cite{jaided2020easyocr} now supports 80+ languages and handles tilted text better out of the box. Ahmed et al.~\cite{cursive_text_recognition_crnn2020} showed that deep CRNN architectures can decode cursive natural-scene text, though at inference costs that are hard to justify on a mobile device. Devi and Rajendran~\cite{ocr_integrated_assistive2023} and Sharma and Gupta~\cite{reading_assistant_blind_ai2022} each built integrated reading assistants for the blind but relied on a single OCR backend, which fails silently on degraded inputs.

VVA's OCR pipeline (\texttt{core/ocr/\_\_init\_\_.py}) takes a different approach: three-tier fallback (EasyOCR $\to$ Tesseract $\to$ MSER region detector), with CLAHE contrast enhancement, bilateral denoising, and Hough-line deskew applied before the first engine is invoked. IOU-based deduplication merges overlapping text regions across backends. During our internal testing we found that the dual-backend fallback recovered roughly 12\% more words on low-contrast signage photographs than either engine alone---though this number should be validated on a standard benchmark.

\subsection{Braille Recognition}

Li et al.~\cite{li2020optical_braille} framed optical braille recognition as semantic segmentation, achieving 98.7\% character accuracy under controlled lighting. Patil and Jadhav~\cite{braille_tts_conversion2021} focused on Braille-to-speech but used a fixed grid detector that failed on embossed paper with variable dot spacing. VVA's braille module (\texttt{core/braille/}) separates plate detection, dot segmentation, character classification, and embossing guidance into four pipeline stages, allowing each to be swapped or retrained independently.

\subsection{Navigation and Spatial Awareness}

Assistive navigation systems have evolved from GPS-only wayfinding to sensor-fused, obstacle-aware guidance. Real and Araujo~\cite{sensors_wearable_navigation2018} surveyed wearable navigation systems and identified latency, battery life, and social acceptability as the three highest-priority user requirements---all of which inform VVA's design. Guerreiro et al.~\cite{outdoor_navigation_vi2022} showed that even with good localization, users need sequential verbal descriptions that build a mental model of the environment. Lee et al.'s StereoPilot~\cite{stereopilot_wearable2023} mapped stereo depth to 3D audio for target localization, but the reliance on stereo cameras limits portability.

More recent work combines multimodal sensors: Rashid et al.~\cite{multimodality_obstacle_nav2023} fused ultrasonic, LIDAR, and camera feeds for obstacle detection, while Saha et al.~\cite{infrastructure_guided_nav2021} anchored navigation to RFID-tagged infrastructure in smart buildings. These approaches achieve high accuracy but presuppose environmental instrumentation. Work on crosswalk scene understanding~\cite{dynamic_crosswalk2022} and shopping complex navigation~\cite{supporting_navigation_shopping2022} further illustrates the domain-specific tuning required. VVA intentionally depends only on a single monocular camera and microphone to work in unmodified environments.

A notable contradiction appears in the literature on spatial audio for navigation: StereoPilot~\cite{stereopilot_wearable2023} reports that users learned 3D audio cues within minutes, whereas Farcy et al.~\cite{tools_tech_blind_nav_review2023} caution that cognitive load from continuous spatial audio leads to fatigue within 30 minutes on average. VVA side-steps this tension by defaulting to natural-language micro-navigation cues (``Stop! Chair half meter left'') rather than spatial audio, with an option for verbose mode when the user explicitly requests richer descriptions.

\subsection{Multimodal Fusion and Scene Understanding}

The convergence of vision, language, and audio has produced powerful scene-understanding models. Radford et al.'s CLIP~\cite{radford2021clip} enabled zero-shot visual recognition from text prompts; Li et al.'s BLIP-2~\cite{li2023blip2} extended this with frozen large language models for open-ended VQA. Antol et al.~\cite{antol2015vqa} defined the VQA benchmark that catalyzed much of this progress. Nguyen et al.~\cite{sbvqa2023} demonstrated speech-based VQA, closing the loop from voice query to visual answer.

In the assistive domain, Park et al.~\cite{emotion_recognition_wearable2022} explored emotion recognition via glasses-type wearables, while Garcia et al.~\cite{enhancing_accessibility_ml_review2024} reviewed machine-learning approaches for both visual and hearing impairments and noted that most systems address a single modality. Indoor scene understanding~\cite{ria_scene_understanding2022} and personal object recognition~\cite{frobt_assistive_robot2019} for visually impaired users have also been explored, though typically in isolation. VVA bridges this gap by fusing 12 distinct capabilities (detection, depth, segmentation, OCR, braille, QR/AR, face recognition, audio localization, action recognition, VQA, memory, and TTS) into a single spoken interface orchestrated by an intent-routing voice controller (\texttt{core/speech/voice\_router.py}).

\subsection{Edge AI and Latency-Constrained Deployment}

Deploying deep networks on consumer hardware remains an active frontier. ONNX Runtime~\cite{onnxruntime2019} provides a portable inference engine, and quantization/pruning pipelines now produce models that fit in a few megabytes. Chen et al.~\cite{eswa_edge_ai_pipeline2020} studied real-time detection on edge devices and showed that INT8 quantization cuts latency by 40\% with less than 1\% mAP drop. LiveKit~\cite{livekit2023} provides WebRTC-based infrastructure that enables VVA to stream video and audio bidirectionally with sub-second round-trip. 

A gap we noticed across the edge-AI literature is evaluating the \emph{full pipeline} cost---not just a single model's latency, but the cumulative cost of detection plus depth plus segmentation plus LLM reasoning plus TTS. VVA's pipeline monitor (\texttt{application/pipelines/pipeline\_monitor.py}) tracks per-stage timings at runtime and triggers fallback paths when cumulative latency exceeds 300\,ms.

\subsection{Summary of Gaps}

Across these clusters, we identify four key gaps: (1)~most assistive systems evaluate single modalities in isolation rather than measuring system-level latency; (2)~few combine more than three sensory modalities into a unified spoken interface; (3)~privacy-preserving on-device memory for personalization is rarely addressed; and (4)~there is no open-source reference implementation with layered architecture enforcement and 400+ tests. VVA aims to address all four.
